\chapter{Verwandte Arbeiten}
Im dritten Kapitel wird der aktuelle Stand der Forschung zum Thema RTOS-Benchmarking dargestellt und diese Arbeit darin eingeordnet. Zuerst wird erläutert, welche Herausforderungen sich beim Vergleich und der Bewertung von Echtzeitsystemen ergeben. Darunter werden auch Fragen zur Reproduzierbarkeit und Vergleichbarkeit zu anderen Betriebssystemen wie FreeRTOS, RIOT oder Contiki-NG bzgl. Speicherverbrauch und Laufzeitverhalten untersucht. Danach werden Zephyr-spezifische Arbeiten behandelt, darunter Zephyrs eigene Referenz-Benchmarks und Studien, die Zephyr in bestimmten Anwendungsszenarien, wie bspw. die Absicherung durch TLS/ DTLS oder KI-Inferenz, einsetzen und teilweise messen. Zuletzt wird diese Arbeit von den vorgestellten Studien abgegrenzt, indem analysiert wird, worin sich ihre Methodik, also der Vergleich von Konfigurationsstufen innerhalb desselben RTOS auf einer spezifizierten Zielplattform, von den anderen bestehenden Vergleichsstudien unterscheidet.

\section{RTOS-Benchmarking allgemein}
Das Benchmarking von Echtzeitbetriebssystemen besteht in der Literatur aus einer Kombination von Laufzeit- und Speichermetriken. Silva et al. (2019) definieren für den Vergleich von Contiki-NG, RIOT und Zephyr vier Bewertungsdimensionen, darunter Leistung, Zeit-Determinismus, Speicher-Fußabdruck und Energieverbrauch. Für die Leistungs-Bewertung benutzen Silva et al. die Thread-Metric Benchmark Suite. Sie besteht aus acht Tests, darunter kooperative und präemptive Kontextwechsel, Interrupt-Verarbeitung, Semaphor-Operationen, Speicherallokation und -freigabe (Silva et al., 2019, S. 10379). Die Zeitvorhersagbarkeit untersuchen die Autoren zudem anhand eigener Mikrobenchmarks für ausgewählte Thread-Verwaltungsfunktionen. Darüber hinaus beziehen sie einen Thread-Metric-Test ein (Silva et al., 2019, S. 10380). Diese Sammlung an Benchmark-Kategorien ähnelt den in dieser Arbeit gewählten fünf Metriken. Hierbei wurde in dieser Arbeit jedoch nicht die Thread-Metric Benchmarks-Suite verwendet, was genauer in Abschnitt 3.4 begründet wird.

Für die Messung des Zeit-Determinismus benutzen Silva et al. (2019, S. 10379) eigenentwickelte Mikrobenchmarks für vier Thread-Verwaltungs-APIs, darunter „Thread Create“, „Resume“, „Suspend“ und „Delete“. Dabei protokollieren sie Minimum, Maximum und Mittelwert als Jitter-Indikator. Auch Leotescu und Vlad (2018) benutzen Mikrobenchmarks in ihrem Ansatz. Sie entwickeln mit „heap\_bench“, „thread\_bench“, „mutex\_bench“ und „context\_switch“ vier eigene, in C geschriebene Testprogramme, die jeweils eine einzelne Kernel-Operation isoliert messen (Leotescu \& Vlad, 2018, S. 3-7). Methodisch ähnelt dies hierbei dem Aufbau dieser Arbeit, wobei hier ebenfalls jede Metrik in einer eigenen, isolierten \lstinline[style=fileinline]|bm_*.c|-Datei gemessen wird (siehe Kapitel 4.1). Außerdem benutzen Leotescu und Vlad (2018, S. 3) für die Zeitmessungen unter Zephyr bewusst die Funktion \lstinline[style=fileinline]|TIMING_INFO_PRE_READ()| anstatt der POSIX-Funktion \lstinline[style=fileinline]|clock_gettime()|, mit der Begründung, dass die POSIX-Funktion aufgrund des Zephyr-Schedulers inkonsistente Werte liefere. Auch diese Arbeit benutzt mit \lstinline[style=fileinline]|timing_counter_get()| eine Zephyr-spezifische Zeitmessschnittstelle anstatt einer POSIX-Zeitabfrage. Dies entspricht der allgemeinen Entscheidung von Leotescu und Vlad für eine hardwarenahe Zephyr-Zeitbasis. Ihr dabei verwendetes Makro \lstinline[style=fileinline]|TIMING_INFO_PRE_READ()| soll jedoch mit der hier eingesetzten API unterschieden werden (Leotescu \& Vlad, 2018, S. 3).

Jaskani et al. (2019) ergänzen dabei diese beiden messtechnisch orientierten Arbeiten um eine breitere, qualitative Einordnung. Sie erwähnen nämlich sieben Auswahlkriterien für ein IoT-Betriebssystem, darunter Architektur, Programmiermodell, Scheduling-Strategie, Speicher-Fußabdruck, Vernetzung, Portabilität und Energieeffizienz (Jaskani et al., 2019, S. 3-4). Zudem ordnen sie Zephyr darin als Mikrokernel-/ Nanokernel-Hybrid mit präemptiven, prioritätsbasierten und EDF-fähigem Scheduler ein (Jaskani et al., 2019, S. 5). Diese Einordnung wird jedoch in dieser Arbeit nicht selbst empirisch nachgeprüft, da der Fokus hier nicht auf einem Architektur-, sondern Konfigurationsvergleich liegt.

\section{Bestehende RTOS-Benchmarks}
Silva et al. benutzen für die betreffende Leistungs-Auswertung 1000 Stichproben. Damit liegt der Stichprobenumfang in derselben Größenordnung wie bei den zeitbezogenen Benchmarks dieser Arbeit, ohne dass daraus schon eine Gleichheit der Messbedingungen folgt (Silva et al., 2019, S. 10379).

Für den untersuchten Zephyr-Aufbau berichten die Autoren außerdem einen mittleren RAM-Bedarf von 52,9 KB und einen Flash-Bedarf von 130,5 KB. Die Werte für Contiki-NG betragen 29,8 KB bzw. 50,6 KB und für RIOT 33,0 KB bzw. 59,9 KB. Diese Zahlen basieren dabei auf die jeweils untersuchten Firmware-Aufbauten inklusive Anwendung und Netzwerkfunktionen und nicht auf einen allgemeinen Mindestbedarf der Betriebssysteme (Silva et al., 2019, S. 10381).

Bei den untersuchten Leistungs- und Zeitvorhersagbarkeits-Metriken dokumentieren Silva et al. schlechtere Ergebnisse für Zephyr als für die beiden Vergleichssysteme. Sie diskutieren dabei, dass weitere Verarbeitungs- und Moduswechsel eine mögliche Erklärung sein können. Auf der anderen Seite notieren sie jedoch, dass der Zeitstempel in ihrem Messaufbau nicht immer am vorgesehenen Endpunkt erfasst wird. Die beobachteten Schwankungen dürfen daher nicht unbegrenzt als den realen Jitter der untersuchten API-Operationen interpretiert werden (Silva et al., 2019, S. 10381).

Für diese Arbeit ist daraus vor allem die Bedeutung von einer klaren Definition und Überprüfung der Messpunkte relevant. Die Ergebnisse von Silva et al. dienen als Vergleichskontext, ersetzen aber dabei keine Validierung der Zeitmessung, die hier eingesetzt werden.

Für diese Arbeit ist diese Erkenntnis auf zwei Arten relevant. Erstens zeigt er, dass Zephyrs Fußabdruck und Determinismus-Verhalten stark von der aktivierten Konfiguration abhängen können. Genau dieser Zusammenhang wird in dieser Arbeit durch die drei Konfigurationsstufen untersucht. Silva et al. beschreiben die Bedingungen ihres Versuchsaufbaus und benutzen dabei Standard- bzw. empfohlene Einstellungen. Eine komplette Aufstellung aller eingeschalteten Zephyr-Konfigurationssymbole wird jedoch nicht angegeben. Zudem werden innerhalb der Untersuchung unter anderem Thread- und Netzwerkbedingungen variiert. Der Vergleich ist deshalb nicht mit einem gezielten Vergleich von den drei hier definierten Konfigurationsprofile gleichzusetzen (Silva et al., 2019, S. 10378–10380).

Die veröffentlichten Werte machen eine Einordnung der Größenordnung möglich. Unterschiede in Hardware, Softwarestand, Anwendung und Messverfahren begrenzen dabei jedoch den direkten Vergleich der Zahlen. Daher bestätigt eine ähnliche Größenordnung allein nicht direkt die Richtigkeit der eigenen Messmethode.

Altoumi und Ahmed beziehen sich dabei UL Solutions Kontextwechselwerte von 223 CPU-Zyklen für FreeRTOS und 524 CPU-Zyklen für Zephyr. Diese Zahlen sind in diesem Kontext keine eigenen Messwerte der Autoren. Sie diskutieren außerdem die Bedeutung von unterschiedlichen MPU- bzw. Sicherheitseinstellungen für die Vergleichbarkeit von solchen Ergebnissen (UL Solutions, 2022, zitiert nach Altoumi \& Ahmed, 2025, S. 2–3; Altoumi \& Ahmed, 2025, S. 3).

Für diese Arbeit ist daran vor allem der methodische Hinweis relevant, dass Schutzfunktionen bei Zeitmessungen genau dokumentiert werden sollen. Die Minimal-Konfiguration deaktiviert die MPU. Dadurch beziehen sich die Werte, die ermittelt werden, auf eine reduzierte Schutzkonfiguration. Ohne einen weiteren vergleichbaren Messlauf mit aktivierter MPU kann sich der vermiedene Aufwand, der durch diese Abschaltung entsteht, nicht getrennt quantifizieren.

Jaskani et al. (2019) dient in dieser Arbeit an erster Stelle als Kontext- und Einordnungsquelle, nicht als Messreferenz. Das liegt daran, dass das Paper selbst keine eigenen Benchmarks durchführt. Stattdessen fasst sie bestehende Literatur zu sechs IoT-Betriebssystemen, darunter Contiki, TinyOS, RIOT, Zephyr, Mbed und Brillo zusammen (Jaskani et al., 2019, S. 5-6).

\section{Zephyr-spezifische Benchmarks}
Ergänzend zu den in 3.2 vorgestellten RTOS-übergreifenden Vergleichsstudien bringt Zephyr selbst eigene Referenz-Benchmarks im Quellbaum mit, der unter \lstinline[style=fileinline]|tests/benchmarks/| liegt. Diese Benchmarks liefern Zephyr-interne Referenzmessungen für ausgewählte Kerneloperationen. In dieser Arbeit werden ihre dokumentierten Messkategorien als Nachvollziehbarkeits- und Abgrenzungspunkte benutzt. Eine unabhängige Validierung der eigenen Messmethodik entsteht daraus allein jedoch noch nicht (Zephyr Project, 2024, tests/benchmarks/latency\_measure/README.rst). Im Folgenden wird für jede der fünf untersuchten Metriken notiert, ob ein Zephyr-Äquivalent existiert, was davon übernommen und was bewusst anders gemacht worden ist. Ergänzend dazu werden drei Studien erwähnt, die Zephyr nicht mit anderen RTOS vergleichen, sondern eine Zephyr-Anwendung auf realer Hardware messen.

\textbf{Kontextwechsel-Latenz }Für diese Metrik existiert durch \lstinline[style=fileinline]|tests/benchmarks/latency_measure| ein Zephyr-eigenes Äquivalent, das laut der darin enthaltenen \lstinline[style=fileinline]|README|-Datei unter anderem die Kontextwechselzeit über \lstinline[style=fileinline]|k_yield()| misst, sowohl präemptiv als auch kooperativ. Es wird zudem auch die Kontextwechselzeit durch Semaphore-Weckung gemessen. Dieser Benchmark wurde in 5.1 schon als Vergleichspunkt genannt. Wichtig ist hierbei jedoch eine Einschränkung. Der genaue Quellcode von \lstinline[style=fileinline]|latency_measure| wurde im Rahmen dieser Arbeit nicht Zeile für Zeile aus dem GitHub-Repository übernommen, sondern nur über die Dokumentation durch die \lstinline[style=fileinline]|README|-Datei eingeordnet. Aus der README werden nur die dort beschriebenen Messkategorien übernommen. Die eigene Kalibrierung des Messaufwands ist davon unabhängig entwickelt worden. Die eigene Auswertung führt Minimum, Maximum, Mittelwert und Stichproben-Standardabweichung laufend über N=1000 Wiederholungen. Sie speichert dabei jedoch keine komplette Verteilung der Einzelwerte.

\textbf{Interrupt-Latenz }Auch dies ist in Zephyrs \lstinline[style=fileinline]|tests/benchmarks/latency_measure| vorhanden. Zephyrs Suite erfasst nämlich die Zeit vom Ende der ISR bis zur Fortsetzung eines Threads. Der in dieser Arbeit gemessene Schritt davor, also vom Software-Trigger bis zur ersten ISR, wird von \lstinline[style=fileinline]|latency_measure| jedoch nicht erfasst. Diese Arbeit misst eine Zeitgröße derselben Größenordnung, wobei der Hinweg vom Auslösen bis zur ersten ISR-Instruktion statt des Rückwegs gemessen wird. Die README beschreibt für \lstinline[style=fileinline]|latency_measure| den Rückweg von einer ISR zum unterbrochenen oder zu einem anderen Thread. Diese Arbeit untersucht auf der anderen Seite den Hinweg vom softwareseitigen Setzen einer NVIC-Anforderung bis zum Zeitstempel am Beginn des eigenen C-Handlers. Aussagen über die bestimmte interne Triggerquelle von \lstinline[style=fileinline]|latency_measure| werden daraus nicht abgeleitet.

\textbf{Timer-Jitter }Im für diese Arbeit geprüften Benchmark-Verzeichnis von Zephyr v3.7 wurde kein Referenzbenchmark gefunden, der die Abweichung der Callback-Intervalle eines periodischen \lstinline[style=fileinline]|k_timer| von seiner Sollperiode misst (Zephyr Project, 2024, tests/benchmarks/).

\textbf{Heap-Allokationsgeschwindigkeit und -Fragmentierung }Für Kernel-Operationen gibt es durch \lstinline[style=fileinline]|tests/benchmarks/sys_kernel| ein Zephyr-eigener Benchmark. Dieser misst jedoch nur Semaphore, LIFO, FIFO, Stack und Memory Slab durch \lstinline[style=fileinline]|k_mem_slab_alloc| und \lstinline[style=fileinline]|k_mem_slab_free|. Es wird also nicht der allgemeine System-Heap über \lstinline[style=fileinline]|k_malloc()| und \lstinline[style=fileinline]|k_free()| untersucht. Ein passenderes Zephyr-Äquivalent wäre dafür \lstinline[style=fileinline]|latency_measure| selbst, da es durch \lstinline[style=fileinline]|heap.malloc.immediate| und \lstinline[style=fileinline]|heap.free.immediate| die Heap-Allokation durchführt. Dabei liefert es, laut der README, in der Standardkonfiguration von Zephyr v3.7.0 Referenzwerte von durchschnittlich 627 Zyklen für \lstinline[style=fileinline]|k_malloc()| und 432 Zyklen für \lstinline[style=fileinline]|k_free()|. Für die Heap-Fragmentierung als Zustandsmaß ist kein Zephyr-eigener Benchmark gefunden worden, da \lstinline[style=fileinline]|sys_heap| selbst keine öffentliche Kennzahl für den größten zusammenhängenden Freiblock bringt (siehe 5.5). Dies ist auch der Grund für die Eigenentwicklung der Bisektions-Probe in dieser Arbeit.

\textbf{Zephyr in Anwendungsszenarien }Leotescu und Vlad untersuchen vier Mikrobenchmarks für Speicher-, Thread-, Mutex- und Kontextwechsel-Operationen. Dabei vergleichen sie Zephyr auf einem i.MX RT1050 mit Cortex-M7 bei 600 MHz mit Linux auf einem i.MX 6UL mit Cortex-A7 bei 528 MHz. Die verschiedene Hardware begrenzt den direkten Vergleich der Betriebssysteme (Leotescu \& Vlad, 2018, S. 1–3).

Für die Zeitmessung unter Zephyr benutzen die Autoren \lstinline[style=fileinline]|TIMING_INFO_PRE_READ()| anstatt \lstinline[style=fileinline]|clock_gettime()|, weil sie mit der POSIX-Zeitabfrage inkonsistente Messwerte beobachtet haben. Dies unterstreicht hierbei die Bedeutung der ausgewählten Zeitbasis. Auch diese Arbeit benutzt für ihre kurzen Laufzeitmessungen eine Zephyr-spezifische Timing-Schnittstelle. Die Eignung muss jedoch für den eigenen Messaufbau beurteilt werden (Leotescu \& Vlad, 2018, S. 3).

In drei wiederholten Benchmark-Läufen notieren die Autoren gleiche zusammengefasste Ergebnisse unter Zephyr. Für den Kontextwechsel nennen sie einen Mittelwert von 47 Zyklen. Unter Linux beobachten sie Abweichungen von bis zu neun Prozent vom jeweiligen Mittelwert. Diese Ergebnisse gelten dabei für den beschriebenen Versuchsaufbau. Sie zeigen jedoch nicht die Verteilung aller einzelnen Operationen. Außerdem zeigen sie auch nicht, dass unabhängige Wiederholungsläufe durch viele Messungen innerhalb eines einzigen Laufs ersetzt werden können (Leotescu \& Vlad, 2018, S. 8–9).

Lee et al. untersuchen TLS-/ DTLS-Client-Funktionalität auf Zephyr zuerst in QEMU und anschließend auf einem FRDM-K64F-Board. Das Board benutzt einen Cortex-M4 und verfügt über 1 MB Flash und 256 KB RAM. Bei den Hardwaretests notieren die Autoren Kommunikationsprobleme und Paketverluste, die dabei in den vorherigen Simulationstests nicht sichtbar waren (Lee et al., 2018, S. 1293–1294).

Die Arbeit dient hier als Beispiel dafür, dass ein erfolgreicher Simulationstest nicht direkt alle Bedingungen von einem realen Gerät abbildet. Sie ist dadurch eine Kontextquelle für die Hardware der IoT-Anwendung, aber keine direkte Referenzmessung für die untersuchten Kernel-Latenzen.

Zhang et al. portieren Paddle Lite auf Zephyr und vergleichen die Inferenz auf einem RK3568 mit Cortex-A55 mit einer Linux-Ausführung. Für die untersuchten Modelle notieren sie dabei durchschnittlich sieben Prozent geringere Inferenzzeiten und einen um 78 Prozent geringeren Systemspeicherverbrauch während der Inferenz unter Zephyr (Zhang et al., 2026, S. 263–264).

Bei diesen Ergebnissen handelt es sich um einen Anwendungsprozessor und eine KI-Inferenzlast. Sie sind daher nicht direkt auf die Cortex-M4-Plattform und die Kernel-Mikrobenchmarks dieser Arbeit übertragbar. Die Studie zeigt jedoch dafür, dass Zephyr auch für größere Anwendungen als ressourcenschonende Ausführungsumgebung untersucht wird.

\section{Abgrenzung der eigenen Arbeit}
Durch die vorgestellten Arbeiten lässt sich die Position dieser Arbeit ableiten. Vier Aspekte unterscheiden sie von den genannten Studien.

Erstens vergleicht diese Arbeit nicht, wie die meisten der vorgestellten Studien (Silva et al., 2019; Altoumi \& Ahmed, 2025; Jaskani et al., 2019), unterschiedliche RTOS-Punkte miteinander. Stattdessen wird ein einzelnes RTOS, nämlich Zephyr, über drei Konfigurationsstufen hinweg, darunter Minimal, Logging und IoT-Stack, verglichen. Diese Methodik beantwortet daher keinen Produktvergleich wie „Welches RTOS ist am schnellsten?“, sondern den Einfluss von Konfigurationsentscheidungen auf Laufzeit- und Speicherverhalten desselben Systems. Silva et al. (2019) und Altoumi und Ahmed (2025) bringen zwar interessante Vergleichswerte zwischen RTOS, sagen aber wenig darüber aus, welcher Teil der beobachteten Unterscheide durch die Konfiguration entsteht, anstatt durch das RTOS-Produkt. Diese Arbeit untersucht genau diese Lücke, indem sie die Konfiguration zur einzigen unabhängigen Variable macht (siehe 4.1).

Zweitens untersucht diese Arbeit Laufzeitmetriken und auch den statischen RAM-/ROM-Fußabdruck für die eigenen Zephyr-Konfigurationsstufen auf einer festgelegten Zielplattform. Die Kombination dieser Metrik-Arten ist nicht neu. Silva et al. Achten schon auf Leistung, Zeitvorhersagbarkeit, Speicher und Energieverbrauch. Der weitere Schwerpunkt dieser Untersuchung liegt auf der Variation der Konfiguration innerhalb desselben Betriebssystems (Silva et al., 2019, S. 10378–10382).

Drittens legt diese Arbeit vor allem Wert auf die Nachvollziehbarkeit des eigenen Versuchsaufbaus. Dazu werden die verwendeten Softwarestände, Konfigurationsfragmente, Messpunkte, Kalibrierungsschritte und Auswertungsgrößen dokumentiert. Diese Angaben sollen dabei eine Wiederholung der Untersuchung erleichtern. Die daraus resultierenden Ergebnisse sagen jedoch nicht aus, dass die Dokumentation den Vergleichsstudien überlegen ist.

Viertens dienen Beobachtungen aus der Literatur als Ausgangspunkt für eigene Entwurfsentscheidungen. Die Diskussion von unterschiedlichen Schutzkonfigurationen bei Altoumi und Ahmed motiviert bspw., die MPU-Einstellung bewusst zu dokumentieren. Die Ergebnisse zum Speicherbedarf von Silva et al. sind zudem ein Grund für die Untersuchung des Speicherbedarfs von unterschiedlichen Zephyr-Aufbauten. Diese Arbeit übernimmt dadurch Fragestellungen aus der Literatur, ohne deren Experimente komplett zu kopieren (Altoumi \& Ahmed, 2025, S. 2–3; Silva et al., 2019, S. 10381).

Insgesamt ist diese Arbeit als eine vertiefte Einzelfallanalyse eines RTOS und kontrollierter Konfigurationsvariation positioniert, die bestehende Literatur-Beobachtungen behandelt und diese auf einer einzigen, genau spezifizierten Hardware-Plattform empirisch nachprüft.
